%%%%%%%%%%  scan idx 345  |  folio 332  %%%%%%%%%%
\SGAPrintPage{332}%
\sgaStableTarget{sga6.e05.prop.3-11}\textbf{Proposition 3.11.} \uline{Let} \(A\) \uline{be a} \(\lambda\)-ring, \uline{and} \(y_i\), \(i=1,\dots,n\) \uline{elements of} \(A\) \uline{such that} \(\lambda^j(y_i)=0\) \uline{for} \(j\ge2\). \uline{Put} \(x=\sum_i y_i\); \uline{then}:
\begin{enumerate}[label=\roman*)]
\item \uline{for} \(k>n\), we have \(\gamma^k(x-n)=0\)
\item \uline{for} \(1\le k\le n\), \(\gamma^k(x-n)\) \uline{is the} \(k\)-th \uline{symmetric function of the} \(y_i-1\).
\end{enumerate}
Let \(j\ge2\). We then have \(\gamma^j(y_i-1)=\lambda^j(y_i+j-2)=\sum_{p+q=j}\lambda^p(y_i)\lambda^q(j-2)\). Since \(\lambda^p(y_i)=0\) for \(p\ge2\), and \(\lambda^q(j-2)=0\) for \(q>j-2\), \(\gamma^j(y_i-1)=0\). The proposition then follows from \(\gamma_t(x-n)=\prod_{i=1}^{n}\gamma_t(y_i-1)\).

\sgaStableTarget{sga6.e05.sec.4}\section*{4. \(\lambda\)-rings generated by generators and relations}

We still work over \(\mathcal C_0\).

\sgaStableTarget{sga6.e05.para.4-1}4.1. Let \(K\) be a \(\lambda\)-ring. Consider a family of indeterminates \((X_i)_{i\in I}\), and let \(A=K[X_i,\lambda^n X_i]_{i\in I,\ n\ge2}\); we wish to define on \(A\) a \(\lambda\)-ring structure extending that of \(K\).

To specify a homomorphism \(\lambda_t:A\to\Ghat(A)\), it suffices to specify the images of the \(X_i\), \(\lambda^n X_i\), which defines a ring homomorphism: \(A\to\Ghat(A)_o\); for it to define a \(\lambda\)-ring structure on \(A\), we must verify the commutativity of diagram (3.7.1). We therefore put:
\sgaStableTarget{sga6.e05.eq.4-1-1}\begin{equation}\tag{4.1.1}
  \lambda_t(X_i)=1+X_it+\sum_{n\ge2}\lambda^n X_i\,t^n~;
  \qquad
  \lambda_t(\lambda^n X_i)=\lambda^n(\lambda_t(X_i)).
\end{equation}
Since diagram (3.7.1) is a diagram of ring homomorphisms, it suffices to verify its commutativity on the \(X_i\), \(\lambda^n X_i\) (it is known by hypothesis on the elements of \(K\)); this follows readily from definitions (4.1.1) and the fact that \(\lambda_u\) and \(\Ghat^u(\lambda_t)\) are \(\lambda\)-homomorphisms.

%%%%%%%%%%  scan idx 346  |  folio 333  %%%%%%%%%%
\SGAPrintPage{333}%
The \(\lambda\)-ring \(A\) is characterized by the following universal property:

\sgaStableTarget{sga6.e05.prop.4-2}\textbf{Proposition 4.2.} \uline{For every} \(K\)-\(\lambda\)-\uline{algebra} \(B\), \uline{and every family} \((b_i)_{i\in I}\) \uline{of elements of} \(B\), \uline{there exists exactly one} \(K\)-\(\lambda\)-\uline{homomorphism} \(f\) \uline{from} \(A\) \uline{to} \(B\) \uline{such that for every} \(i\in I\), \(f(X_i)=b_i\).

For \(f\) to be a \(\lambda\)-homomorphism, we must have \(f(\lambda^j X_i)=\lambda^j(b_i)\) for all \(i,j\); this defines \(f\) uniquely. The homomorphism thus obtained is indeed a \(\lambda\)-homomorphism, since to verify that \(\lambda_t\circ f=\Ghat(f)\circ\lambda_t\) it suffices to verify it on a system of generators of \(A\) over \(K\), which follows from the definition of \(f\) and relations (2.4.2).

\sgaStableTarget{sga6.e05.para.4-3}4.3. We can make the same construction using the \(\gamma\) in place of the \(\lambda\). Let \(A'=K[X_i,\gamma^n X_i]_{i\in I,\ n\ge2}\); to define the \(\lambda\)-ring structure on \(A'\), we specify a ring homomorphism: \(A'\to\Ghat(A')_\nabla\), by:
\sgaStableTarget{sga6.e05.eq.4-3-1}\begin{equation}\tag{4.3.1}
  \gamma_t(X_i)=1+X_it+\sum_{n\ge2}\gamma^n X_i\mathbin{.}t^n~;
  \qquad
  \gamma_t(\gamma^n X_i)=\gamma'^n(\gamma_t(X_i)).
\end{equation}
We then verify the commutativity of (3.7.2), which shows that we do obtain a \(\lambda\)-ring structure.

The \(\lambda\)-ring \(A'\) also satisfies the universal property 4.2; it is therefore canonically isomorphic (as a \(\lambda\)-ring) to the \(\lambda\)-ring \(A\).

\sgaStableTarget{sga6.e05.def.4-4}\textbf{Definition 4.4.} \uline{The} \(\lambda\)-\uline{ring} \(A\) \uline{will be called the} free \(\lambda\)-ring generated over \(K\) by the family of generators \((X_i)_{i\in I}\).

\sgaStableTarget{sga6.e05.def.4-5}\textbf{Definition 4.5.} \uline{Let} \(A\) \uline{be a} \(\lambda\)-\uline{ring}, \(J\) \uline{an ideal of} \(A\). \uline{We say that} \(J\) \uline{is} \uline{a} \(\lambda\)-ideal of \(A\) \uline{if for every} \(x\in J\), \uline{and every} \(n\ge1\), \(\lambda^n(x)\in J\).

If \(f:A\to B\) is a \(\lambda\)-homomorphism, the kernel of \(f\) is a \(\lambda\)-ideal.

Let \(J\) be a \(\lambda\)-ideal; the composite homomorphism: \(A\xrightarrow{\ \lambda_t\ }\Ghat(A)\longrightarrow\Ghat(A/J)\) vanishes on \(J\) and therefore factors through \(A/J\), which is thus endowed
%%%%%%%%%%  scan idx 347  |  folio 334  %%%%%%%%%%
\SGAPrintPage{334}%
with a \(\lambda\)-ring structure, relations (2.4.2) passing to the quotient.

Every intersection of \(\lambda\)-ideals is a \(\lambda\)-ideal. In particular, there exists a smallest \(\lambda\)-ideal containing a given family of elements \((x_i)_{i\in I}\); by relations (2.4.2), it is the (ordinary) ideal generated by the \(\lambda^n(x_i)\), \((n\ge1)\).

Again let \(K\) be a \(\lambda\)-ring, \((X_i)_{i\in I}\) a family of indeterminates, \(A=K[(X_i),(\lambda^n X_i)]\), and let \((P_{i'})_{i'\in I'}\) be a family of elements of \(A\): let \(J\) be the \(\lambda\)-ideal generated by the \(P_{i'}\).

\sgaStableTarget{sga6.e05.def.4-6}\textbf{Definition 4.6.} \uline{The} \(\lambda\)-\uline{ring} \(A/J\) \uline{will be called the} \(\lambda\)-ring generated over \(K\) by the family of generators \((X_i)_{i\in I}\) subject to the relations \((P_{i'})_{i'\in I'}\).

We can also specify relations between the \(X_i\), \(\gamma^n X_i\) and make the same construction, observing that an ideal is a \(\lambda\)-ideal if and only if it is stable under the operations \(\gamma^n\), \(n\ge1\), as follows from relations 3.2.

Let \(B\) be a \(K\)-\(\lambda\)-algebra, \((b_i)_{i\in I}\) a family of elements of \(B\). Since the \(P_{i'}\) are polynomials in the indeterminates \(X_i\), \(\lambda^n X_i\), we can substitute the \(b_i\), \(\lambda^n(b_i)\) for these indeterminates; if the resulting element is zero, we say that the \(b_i\) satisfy the relations \(P_{i'}\). The \(K\)-\(\lambda\)-algebra \(A/J\) is then characterized by the following universal property:

\sgaStableTarget{sga6.e05.prop.4-7}\textbf{Proposition 4.7.} \uline{For every} \(K\)-\(\lambda\)-\uline{algebra} \(B\) \uline{and every family} \((b_i)_{i\in I}\) \uline{of elements of} \(B\), \uline{satisfying the relations} \((P_{i'})_{i'\in I'}\), \uline{there exists exactly one} \(K\)-\(\lambda\)-\uline{homomorphism} \(f':A/J\to B\) \uline{such that for every} \(i\), we have \(f'(X_i)=b_i\).

Indeed, the kernel of the homomorphism \(f\) defined in 4.2 contains the \(P_{i'}\), hence \(J\) since it is a \(\lambda\)-ideal; \(f\) therefore passes to the quotient and defines
%%%%%%%%%%  scan idx 348  |  folio 335  %%%%%%%%%%
\SGAPrintPage{335}%
\(f'\), which is plainly unique.

\sgaStableTarget{sga6.e05.prop.4-8}\textbf{Proposition 4.8.} \uline{Let} \(n\) \uline{be an integer} \(\ge0\); \uline{in the free} \(\lambda\)-\uline{ring generated} \uline{over} \(K\) \uline{by an indeterminate} \(X\), \uline{the ideal generated by the} \(\lambda^k X\), \(k>n\), \uline{is} \uline{a} \(\lambda\)-\uline{ideal}.

Put \(A=K[X,\lambda^k X]_{k\ge2}\), and let \(J\) be the ideal generated by the \(\lambda^k X\) for \(k>n\); put \(B=A/J\). We must show that the composite homomorphism: \(A\xrightarrow{\ \lambda_t\ }\Ghat(A)\longrightarrow\Ghat(B)\) vanishes on \(J\), i.e. that \(\lambda_t(\lambda^k X)\), \(k>n\), has image \(1\) in \(\Ghat(B)\). Now \(\lambda_t(\lambda^k X)=\lambda^k(\lambda_t(X))\), and \(\Ghat(A)\to\Ghat(B)\) is a \(\lambda\)-homomorphism; it therefore suffices to show that for every ring \(C\) and every series \(\varphi\in\Ghat(C)\) whose coefficients vanish in degrees \(>n\), \(\lambda^k(\varphi)=1\) when \(k>n\). But to prove this last property it suffices to prove it for \(C=K[X_1,\ldots,X_n]\), and \(\varphi=1+X_1t+\ldots+X_nt^n\). Embedding \(K[X_1,\ldots,X_n]\) in \(K[T_1,\ldots,T_n]\) by means of symmetric functions, we can write: \(1+X_1t+\ldots+X_nt^n=\prod_{i=1}^{n}(1+T_it)\), hence
\[
  \lambda_u(1+X_1t+\ldots+X_nt^n)=\prod_{i=1}^{n}\lambda_u(1+T_it)=\prod_{i=1}^{n}\bigl(\underline{1}+\{1+T_it\}u\bigr),
\]
which has degree \(n\) in \(u\).

\sgaStableTarget{sga6.e05.cor.4-9}\textbf{Corollary 4.9.} \uline{The} \(\lambda\)-\uline{ring generated over} \(K\) \uline{by the generator} \(X\) \uline{subject to the relations} \(\lambda^k(X)=0\) \uline{for} \(k>n\), \uline{is isomorphic to} \(K[X,\ldots,\lambda^n X]\).

4.9.1. We would have the same statements on replacing the \(\lambda\) by the \(\gamma\). We can also replace \(X\) by a family of indeterminates \((X_i)_{i\in I}\), subject to the relations \(\lambda^k(X_i)=0\) for \(k>n_i\), where the \(n_i\) are integers \(\ge2\); the resulting \(\lambda\)-ring is then \(K[X_i,\lambda^k X_i]_{i\in I,\ 2\le k\le n_i}\).

\sgaStableTarget{sga6.e05.para.4-10}4.10. Let \(K\) be a binomial ring, and \(A\) an augmented \(K\)-\(\lambda\)-algebra. We denote by \(B=A[X,\gamma^k X]_{k\ge2}\) the free \(\lambda\)-ring generated by \(X\) over \(A\);
%%%%%%%%%%  scan idx 349  |  folio 336  %%%%%%%%%%
\SGAPrintPage{336}%
we define a \(\lambda\)-homomorphism: \(B\to K\) extending the augmentation of \(A\) by sending \(X\) to \(0\), which endows \(B\) with an augmented \(K\)-\(\lambda\)-algebra structure. We propose to determine its filtration (3.10) in terms of that of \(A\).

We grade \(B\) by weight with respect to the \(\gamma^k X\), assigning \(\gamma^k X\) the weight \(k\); we denote by \(M_k\) the set of monomials of weight \(k\) in the indeterminates \(\gamma^i X\).

\sgaStableTarget{sga6.e05.lem.4-11}\textbf{Lemma 4.11.} \uline{Let} \(a\in A\); \uline{the coefficient of a monomial of weight} \(n\) \uline{with respect} \uline{to the} \(\gamma^j X\) \uline{in} \(\gamma^k(aX)=P'_k(\gamma^i(a),\gamma^j X)\) \uline{has filtration at least} \(k-n\) (\uline{in} \(A\))

The polynomial \(P'_k\) has no isobaric components of weight \(<k\) in the \(\gamma^i(a)\), \(\gamma^j X\) (3.5), so the lemma is true when \(a\) has zero augmentation.

If \(\epsilon\) denotes the augmentation homomorphism: \(A\to K\), we can write in the general case:
\[
  \gamma_t(aX)=\gamma_t((a-\epsilon(a))X)\mathbin{.}\gamma_t(\epsilon(a)X).
\]
Now \(\gamma_t(\epsilon(a)X)=(\gamma_t(X))^{\epsilon(a)}=\sum_{n=0}^{\infty}R_{n,\epsilon(a)}((\gamma^j X))\mathbin{.}t^n\), the \(R_{n,\epsilon(a)}\) being isobaric of weight \(n\) in the \(\gamma^j X\) (3.8); we deduce the lemma by taking the product.

\sgaStableTarget{sga6.e05.prop.4-12}\textbf{Proposition 4.12.} \sgaStableTarget{sga6.e05.eq.4-12-1}\(\Fil^k B=\Fil^k A+\Fil^{k-1}A\mathbin{.}M_1+\ldots+\Fil^0 A\mathbin{.}M_k+B\mathbin{.}M_{k+1}\).

It is clear that the second ideal is contained in the first; it therefore remains only to prove the reverse inclusion.

Since the filtration defined by the family of ideals on the right-hand side for varying \(k\) is a ring filtration, it suffices to prove the reverse inclusion for an element of the form \(\gamma^k(b)\), with \(\epsilon(b)=0\), these elements
